\documentclass[11pt]{article}
\textwidth 15.9cm
\textheight 23. cm
\addtolength{\oddsidemargin}{-11.5mm}
\addtolength{\evensidemargin}{-11.5mm}
\addtolength{\topmargin}{-19.0mm}

%\usepackage{amsmath}
\usepackage[T1]{fontenc}
\usepackage[ansinew]{inputenc}
\usepackage[bbgreekl]{mathbbol}
\usepackage{amsmath,amssymb}
\usepackage{graphicx}
\usepackage[colorlinks=true,urlcolor=cyan]{hyperref}%,linkcolor=green]{hyperref}
\usepackage{url}
\usepackage[numbers,square]{natbib}
\usepackage{multirow}
\usepackage{tikz}
\usepackage{tikz-cd}
\usetikzlibrary{arrows,matrix}
\usepackage{bbm}
\usepackage{authblk}
\usepackage{subcaption}
\usepackage{xstring}
\usepackage{xparse}
\usepackage{comment}
\renewcommand{\Affilfont}{\footnotesize}

\usepackage{stmaryrd}
\usepackage{comment}
\usepackage{showlabels}
\usepackage{soul}

\usetikzlibrary{calc,patterns,shapes.geometric}
\usetikzlibrary{arrows,shapes,positioning}
\usetikzlibrary{decorations.markings,decorations.pathmorphing,decorations.pathreplacing}
\usetikzlibrary{decorations.text,decorations.shapes}

%%%%%%%% MACROS
\input{GempicLatexMacros}

\title{Coding Style and Conventions}


\author[1]{The Gempic development team}
\affil[1]{Max-Planck-Institut f\"ur Plasmaphysik, Garching, Germany}


\begin{document}

\maketitle

\tableofcontents

\section{Code Editor: VSCode}
You're not required to use any specific editor, but the majority of the development team uses VSCode. Some useful tips in this regard have been gathered below.

\subsection{VSCode Setup}
\begin{enumerate}
   \item Install the extensions: C/C++, C/C++ Extension Pack, Clang-Format, CMake, CMake Tools
   \item Open folder where the cloned gempic is. This will be \verb|$(workspaceFolder)|.
   \item Open settings from the menu (or \verb|Ctrl|+\verb|,| on linux)
   \begin{itemize}
      \item In Extensions/CMake Tools set \verb|Cmake: Source Directory| to \verb|${workspaceFolder}|
      \item Formatting: set \verb|C_Cpp.clang_format_style| to .clang-format (the .clang-format file is in the home of the gempic repository)
   \end{itemize}
\end{enumerate}

\subsection{Configuring, Compiling}
If you have all the required dependencies, GEMPICX should work out of the box using one of the given presets. VSCode automatically detects the presets file and prompts you to select one, but if you want to change your preset, simply open the command prompt from the menu (or \verb|Ctrl|+\verb|Shift|+\verb|P| on linux) and choose \verb|CMake: Select Configure Preset| followed by \verb|CMake: Build|. It is helpful to set the setting \verb|Cmake: Parallel Jobs| for significantly faster build times.

\subsection{Launching executables}
Debugging: Two configuration examples are in scripts/launch.json
\begin{itemize}
   \item vlasov-maxwell: simulation example with input file
   \item (lldb) Cmake: debugs the target that is set in vscode (without input file). Can be used to debug ctests
\end{itemize}

\subsection{Remote work on the compute clusters (Raven/Viper)}
It is possible to work in VSCode over SSH using the extension Remote-SSH. This is an advanced feature in that it takes some fiddling to get to work:
\begin{enumerate}
   \item Set \verb|Remote-SSH: Config File| to point towards your ssh config file.
   \item Check \verb|Remote.SSH: Show Login Terminal|.
   \item Run the command \verb|Remote-SSH: Connect to Host| and follow the steps to add your connection.
\end{enumerate}

An additional few tips might be useful for remote work:
\begin{itemize}
   \item Intellisense works poorly when paths are not absolute, so it's heavily recommended to have your workspace be an absolute path, \emph{e.g.} \verb|/raven/u/YOURUSERNAME/...| instead of \verb|\u\YOURUSERNAME/...|
   \item If you find yourself timing out when you attempt to log in, set \verb|Remote.SSH: Connect Timeout| to a sufficiently high number, such as \verb|30|.
   \item You can set the setting \verb|Cmake: Parallel Jobs| to \verb|2| on a non-interactive login node or \verb|6| on an interactive one.
\end{itemize}

\end{document}
